\chapter{Nitrites (Urine)}

\section*{What Is This Test?}
A urinary nitrite test looks for nitrites in the urine. Nitrites serve as a surrogate marker of bacteriuria. Many Gram-negative uropathogens (Escherichia coli, Klebsiella, Proteus, Enterobacter, Pseudomonas) reduce dietary nitrate to nitrite through bacterial nitrate reductase. Nitrite therefore indicates significant bacterial colonization. The test is a routine dipstick pad. It is most specific (highly predictive when positive) for urinary tract infection when combined with the leukocyte esterase pad. The trade-off is reduced sensitivity: a negative test does not exclude UTI \cite{deville2004urine}. False negatives occur with bladder dwell time $<$ 4 hours, low dietary nitrate, and organisms that do not reduce nitrate (Gram-positive enterococci, Staphylococcus saprophyticus).

\section*{Alternative Names}
\begin{itemize}
  \item Urine Nitrite
  \item Dipstick Nitrite
  \item Nitrituria
  \item Urinalysis Nitrite
\end{itemize}
\section*{Principle}
The dipstick pad uses the Griess reaction. Nitrite reacts with p-arsanilic acid at acidic pH to form a diazonium compound. This couples with N-(1-naphthyl)ethylenediamine to form a pink-red azo dye. Detection threshold: 0.06--0.15 mg/dL nitrite. Color intensity is graded as negative or positive (no scaling, because the test is qualitative) \cite{strasinger2020urinalysis}.

\section*{History}
The Griess reaction was first described in 1858 and adapted to urine dipsticks by the 1950s. Combined with the leukocyte esterase reaction, it became a cornerstone of UTI screening.

\section*{Why This Test Is Ordered}
\begin{itemize}
  \item Screening for urinary tract infection in symptomatic patients (dysuria, frequency, urgency)
  \item Routine urinalysis screening
  \item Evaluation of asymptomatic bacteriuria in pregnancy \cite{nicolle2019clinical}
  \item Pre-operative evaluation before urologic procedures
  \item Long-term care or catheter surveillance (interpret with caution, because colonization is common without infection)
\end{itemize}

\section*{Primary vs Secondary Test}
This is a primary screening test. A positive result with symptoms warrants urine culture confirmation.

\section*{How This Test Is Performed}
A first-morning urine sample is preferred, because it provides the $>$ 4 hours of bladder dwell time that bacteria need to convert nitrate to nitrite. A midstream clean-catch sample is standard. The dipstick pad is read at 60 seconds. A positive result reflexively triggers urine culture and susceptibility testing.

\section*{What Preparation Is Needed}
None specific. An adequate bladder dwell time ($>$ 4 hours) is essential. Patients on low-nitrate diets (vegetarian, restricted diets) may have false negatives. Vitamin C can interfere.

\section*{Contraindications and Precautions}
\begin{itemize}
  \item False negatives: short bladder dwell time ($<$ 4 hours), low dietary nitrate, organisms that do not reduce nitrate (Enterococcus faecalis, Staphylococcus saprophyticus), high urinary vitamin C, frequent voiding, dilute urine. False positives are rare: contaminated specimens, pink-tinted medications (phenazopyridine), and expired dipsticks can cause them. Phenazopyridine (Pyridium) can cause a false positive.
\end{itemize}
\section*{Risks}
None.

\section*{What Happens After the Test}
A positive nitrite with symptoms plus a positive leukocyte esterase strongly suggests UTI: give empirical antibiotics with culture confirmation. A negative nitrite with symptoms does not exclude UTI; proceed with urinalysis microscopy and culture if clinical concern remains.

\section*{Understanding Results}

\subsection*{Below Normal}
Negative nitrite is normal. It does not exclude UTI: Gram-positive organisms, short bladder dwell time, or low nitrate intake all produce false negatives.

\subsection*{Above Normal}
\begin{itemize}
  \item Significant bacteriuria from nitrate-reducing Gram-negative organisms (\textit{E.\ coli} causes 80\% of community UTIs) \cite{hooton2020acute}
  \item Asymptomatic bacteriuria, especially in pregnancy or preoperative evaluation
  \item Catheter-associated asymptomatic bacteriuria in long-term catheterized patients
  \item False positives with medications (phenazopyridine) or stale contaminated specimens
\end{itemize}

\section*{Normal Results}
\begin{itemize}
  \item \textbf{Nitrite (urine dipstick):} Negative in normal urine
  \item \textbf{Interpretation when positive:} Strongly suggests bacteriuria if bladder dwell time is adequate
  \item \textbf{Combined with leukocyte esterase:} Both positive increases specificity for UTI
  \item \textbf{Cutoff concerns:} Detection requires $>$ 10$^5$ CFU/mL in most laboratories
\end{itemize}

\section*{Follow-up Tests}
\begin{itemize}
  \item \textbf{Urine culture and susceptibility:} Required for definitive diagnosis and treatment
  \item \textbf{Urinalysis with microscopic examination:} Confirms leukocyte esterase, WBC, bacteria
  \item \textbf{Leukocyte esterase pad of dipstick:} Indirect detection of pyuria
  \item \textbf{Pregnancy test (if applicable):} If UTI is suspected in women of childbearing age
  \item \textbf{Renal ultrasound:} For recurrent UTIs or suspected obstruction
\end{itemize}

\section*{References}
\bibliographystyle{unsrtnat}
\bibliography{references}

\clearpage
\section*{Regulatory Status and Authorization Date}
This chapter describes a test category, not a specific commercial product. FDA, CDSCO, EU IVDR, and MHRA approval or registration dates therefore cannot be assigned without the assay manufacturer, product name, instrument, and intended use. For laboratory-developed tests, an individual agency approval date may not exist. See the Preface for the interpretation of regulatory status.

\clearpage
